\documentclass[11pt]{article}
\usepackage[margin=1in]{geometry}
\usepackage{parskip}
\usepackage{hyperref}
\usepackage{xcolor}
\usepackage{listings}
\usepackage{booktabs}
\usepackage{graphicx}

\hypersetup{colorlinks=true, linkcolor=blue, urlcolor=blue}
\definecolor{codebg}{rgb}{0.96,0.96,0.96}
\lstset{basicstyle=\ttfamily\small, backgroundcolor=\color{codebg}, breaklines=true,
  columns=fullflexible, frame=single, framesep=3pt, xleftmargin=3pt, xrightmargin=3pt}

\title{Precision Support for Row-Major $\times$ Row-Major GPU Matrix\\
Multiplication: FP16, BF16, and FP32 on Tensor Cores}
\author{ganesh@cs.utah.edu \\ \small (written with Claude Code)}
\date{\today}

\begin{document}
\maketitle

\begin{abstract}
This report studies matrix multiplication $C = A \times B$ where both $A$ and $B$ arrive
in \textbf{row-major} layout, restricted to three precisions: FP16, BF16, and FP32.
Integer (INT8/IMMA) matrix multiplication is out of scope and has been commented out of
\texttt{tc\_matmult/leimao-mm.cu} accordingly. Every claim about what the hardware and
CUDA actually support here is verified directly against the CUDA 12.4
\texttt{mma.hpp} header shipped on this machine (path in
Section~\ref{sec:fp16}), not general recollection --- exact fragment types, tile shapes,
and accumulator types are quoted from it. FP16 is fully exercised, working code
(\texttt{leimao-mm.cu}); BF16 and the FP32/TF32 path are supported by the same header but
not exercised by that file. Section~\ref{sec:simt} answers directly: no SIMT-core
(ordinary per-thread CUDA) matrix multiplication happens anywhere in this file --- every
GPU-side multiply-accumulate goes through tensor cores.
\end{abstract}

\tableofcontents

\section{Scope and What ``Row-Major, Row-Major'' Means Here}
\label{sec:scope}

\texttt{leimao-mm.cu} (Lei Mao's WMMA GEMM example, copied verbatim into this
repository) exposes its GEMM through a column-major-plus-transpose-flag interface:
\begin{lstlisting}
launch_wmma_mm(A, B, C, m, n, k, is_A_transpose, is_B_transpose, stream);
\end{lstlisting}
Its own header comment states the convention precisely: \emph{``All the data in the
matrices are stored in a column-major order... \texttt{WMMA\_FRAG\_LAYOUT\_A}:
\texttt{row\_major} if A is transposed, otherwise \texttt{col\_major}.''} This is a
standard trick, not a real transpose: an $M \times K$ matrix stored \textbf{row-major}
occupies \emph{exactly} the same bytes in memory as its $K \times M$ transpose stored
\textbf{column-major}. So ``\texttt{is\_A\_transpose=true}'' in this code does not mean
``compute $A^T$'' in any mathematical sense that changes the answer --- it means ``the
data \texttt{A} points to is laid out row-major; load it that way.'' The code confirms
this directly: setting the flag selects \texttt{nvcuda::wmma::row\_major} as the fragment
load layout, and adjusts the leading dimension accordingly
(\texttt{lda = is\_A\_transpose ? k : m}, correct for a row-major $M\times K$ array).

\textbf{Consequence for this study:} row-major $A$ and row-major $B$ is obtained with
\begin{lstlisting}
launch_wmma_mm(A, B, C, m, n, k, /*is_A_transpose=*/true, /*is_B_transpose=*/true, stream);
\end{lstlisting}
--- the branch the demo's own \texttt{print\_test\_header} labels ``\texttt{C = A\^{}T *
B\^{}T}''. That label is accurate to the code's internal naming convention but
\emph{not} a mathematical transpose of the matrices this report cares about: it is
exactly $C = A \times B$ with both operands read as ordinary row-major arrays.

\section{FP16 (HMMA): Working, and Already in the File}
\label{sec:fp16}

\subsection{What the header actually defines}
From \texttt{crt/mma.hpp}, CUDA 12.4 (full path on this machine:
\texttt{/usr/local/cuda-12.4/targets/\allowbreak x86\_64-linux/\allowbreak
include/crt/mma.hpp}): \texttt{wmma::fragment<matrix\_a/matrix\_b, 16, 16,
16, \_\_half, row\_major$|$col\_major>} is defined, with \texttt{load\_matrix\_sync} for
both layouts. Two families of \texttt{mma\_sync} overloads exist for the accumulator:
one into \texttt{fragment<accumulator, 16, 16, 16, \_\_half>} (FP16 accumulation) and one
into \texttt{fragment<accumulator, 16, 16, 16, float>} (FP32 accumulation) --- \emph{both}
are real, hardware-backed instruction sequences, not a software emulation of one from the
other. The same three families of tile shapes appear for FP16 as for every other
supported type here: $16\!\times\!16\!\times\!16$, $32\!\times\!8\!\times\!16$, and
$8\!\times\!32\!\times\!16$.

\texttt{leimao-mm.cu} always instantiates \texttt{launch\_wmma\_mm<\_\_half, float>} ---
i.e.\ it always chooses the \textbf{FP32-accumulator} HMMA path, at the
$16\times16\times16$ tile shape, and never exercises the FP16-accumulator variant even
though the hardware/header support it.

\subsection{How the multiplication is realized}
One CUDA \emph{warp} (32 threads acting as one unit) computes one $16\times16$ output
tile per iteration of the outer loop, striding over $K$ in steps of 16:
\begin{lstlisting}
warpM, warpN <- this warp's 2D tile coordinates (from blockIdx/threadIdx)
acc_frag: fragment<accumulator, 16,16,16, float> <- fill_fragment(0)
for ki = 0, 16, 32, ..., k-16:
    a_frag: fragment<matrix_a, 16,16,16, half, row_major>
    b_frag: fragment<matrix_b, 16,16,16, half, row_major>
    load_matrix_sync(a_frag, A + offset(warpM, ki), lda)   # row-major load
    load_matrix_sync(b_frag, B + offset(ki, warpN), ldb)   # row-major load
    mma_sync(acc_frag, a_frag, b_frag, acc_frag)           # one HMMA tensor-core op
c_frag: fragment<accumulator, 16,16,16, float> <- load existing C tile
for each element e in c_frag: c_frag[e] = alpha*acc_frag[e] + beta*c_frag[e]
store_matrix_sync(C + offset(warpM, warpN), c_frag, ldc, mem_col_major)
\end{lstlisting}
Every multiply-accumulate over the $16\times16\times16$ tile happens in a single
\texttt{mma\_sync} call --- one tensor-core HMMA instruction, not 4096 separate
multiply-adds issued by ordinary CUDA cores. \texttt{load\_matrix\_sync} and
\texttt{store\_matrix\_sync} handle the row-major$\leftrightarrow$fragment layout
translation entirely inside the intrinsic; the kernel code never manually transposes
anything.

\section{BF16: Supported by the Header, Not Exercised by This File}
\label{sec:bf16}

\textbf{Yes, BF16 is supported} --- on Ampere (compute capability 8.0+, which includes
\texttt{sm\_86}, the target this file already builds for) and confirmed directly in the
same header:
\begin{lstlisting}
wmma::fragment<matrix_a/matrix_b, 16, 16, 16, __nv_bfloat16, row_major|col_major>
\end{lstlisting}
is defined, at all three tile shapes (\texttt{16$\times$16$\times$16},
\texttt{32$\times$8$\times$16}, \texttt{8$\times$32$\times$16}), with full
\texttt{load\_matrix\_sync} support for both layouts.

\textbf{One structural difference from FP16, confirmed by the header:} only
\texttt{mma\_sync} overloads accumulating into \texttt{fragment<accumulator, 16,16,16,
float>} exist for BF16 --- there is \emph{no} BF16-accumulator overload. BF16 tensor-core
matrix multiplication on this hardware \textbf{always} accumulates in FP32; there is no
choice to make, unlike FP16.

\texttt{leimao-mm.cu} never instantiates a BF16 path --- \texttt{\_\_nv\_bfloat16} does
not appear anywhere in the file. Adding it would be a direct, mechanical change: swap
\texttt{\_\_half} for \texttt{\_\_nv\_bfloat16} as \texttt{T1} (keeping \texttt{T2 =
float}, since that is the only accumulator type available), and replace
\texttt{float2half}'s \texttt{\_\_float2half} conversion with \texttt{\_\_float2bfloat16}.
The kernel body is otherwise identical to Section~\ref{sec:fp16}:
\begin{lstlisting}
acc_frag: fragment<accumulator, 16,16,16, float> <- fill_fragment(0)   # float is mandatory
for ki = 0, 16, 32, ..., k-16:
    a_frag: fragment<matrix_a, 16,16,16, __nv_bfloat16, row_major>
    b_frag: fragment<matrix_b, 16,16,16, __nv_bfloat16, row_major>
    load_matrix_sync(a_frag, A + offset(warpM, ki), lda)
    load_matrix_sync(b_frag, B + offset(ki, warpN), ldb)
    mma_sync(acc_frag, a_frag, b_frag, acc_frag)           # one BF16 tensor-core op
# store identical to FP16 case, c_frag is float throughout
\end{lstlisting}

\section{FP32: Not on Tensor Cores at Full Precision --- Only via TF32}
\label{sec:fp32}

This is the precision where the answer is not ``yes, same pattern, different type.''

\subsection{What the header does \emph{not} define}
Searching \texttt{crt/mma.hpp} for a tensor-core fragment holding true, full-mantissa
\texttt{float} operands (23-bit mantissa, IEEE-754 single precision) for
\texttt{matrix\_a}/\texttt{matrix\_b} finds nothing. \textbf{Tensor cores on this GPU
cannot multiply two full-precision FP32 matrices directly.} There is no WMMA
instruction for it, at any tile shape.

\subsection{What the header does define: TF32}
\texttt{wmma::fragment<matrix\_a/matrix\_b, 16, 16, \textbf{8}, precision::tf32,
row\_major$|$col\_major>} --- note the $K$ dimension is \textbf{8}, not 16, unlike every
other type in this report. \texttt{load\_matrix\_sync} for this fragment still takes a
plain \texttt{const float*} pointer: the caller's data in host/device memory stays
ordinary FP32 the whole time. The precision reduction --- TF32 keeps FP32's 8-bit
exponent but truncates the mantissa to 10 bits, matching FP16's mantissa width while
keeping FP32's exponent range --- happens \emph{inside} the load intrinsic itself, not as
a separate conversion step the caller writes. The accumulator,
\texttt{fragment<accumulator, 16,16,8, float>}, is full-precision \texttt{float}, and
\texttt{mma\_sync} overloads for all four row/col-major combinations of
\texttt{matrix\_a}/\texttt{matrix\_b} are present, confirming TF32 GEMM is fully
supported for row-major operands, exactly like Sections~\ref{sec:fp16}
and~\ref{sec:bf16}.

\texttt{leimao-mm.cu} does not exercise this path either --- no \texttt{precision::tf32}
anywhere in the file. A TF32 version of the kernel:
\begin{lstlisting}
acc_frag: fragment<accumulator, 16,16,8, float> <- fill_fragment(0)
for ki = 0, 8, 16, ..., k-8:                              # K-tile is 8, not 16
    a_frag: fragment<matrix_a, 16,16,8, precision::tf32, row_major>
    b_frag: fragment<matrix_b, 16,16,8, precision::tf32, row_major>
    load_matrix_sync(a_frag, A_float_ptr + offset(warpM, ki), lda)  # FP32 in, TF32 internally
    load_matrix_sync(b_frag, B_float_ptr + offset(ki, warpN), ldb)
    mma_sync(acc_frag, a_frag, b_frag, acc_frag)           # one TF32 tensor-core op
# store: c_frag is float, full precision, no further loss on the output/accumulate side
\end{lstlisting}

\subsection{If true full-precision FP32 GEMM is actually required}
Then it cannot run on tensor cores on this hardware at all --- it has to be an ordinary
SIMT kernel using plain \texttt{FFMA} instructions on the regular CUDA cores, ordinary
per-thread arithmetic, no \texttt{wmma} namespace involved:
\begin{lstlisting}
# one thread per output element (a real kernel would tile through shared memory;
# this is the conceptual/unoptimized form)
row, col <- this thread's output coordinates
acc: float = 0
for ki = 0, 1, 2, ..., k-1:
    acc += A[row, ki] * B[ki, col]        # ordinary FFMA, full 23-bit mantissa throughout
C[row, col] = acc
\end{lstlisting}
\texttt{leimao-mm.cu} does not contain this kernel either (see Section~\ref{sec:simt}) ---
its only full-precision FP32 code is the CPU-side reference function, not a GPU kernel at
all.

\section{Is Any of This SIMT-Core Matrix Multiplication?}
\label{sec:simt}

No. \texttt{leimao-mm.cu} contains exactly one \texttt{\_\_global\_\_} kernel,
\texttt{wmma\_gemm\_a\_col\_major\_b\_col\_major}, and every multiply-accumulate inside
it happens via \texttt{nvcuda::wmma::mma\_sync} --- a tensor-core instruction. There is no
kernel anywhere in this file that computes a matrix product using ordinary per-thread CUDA
arithmetic (\texttt{FFMA}, \texttt{HFMA2}, or otherwise) on the GPU.

The one non-tensor-core matrix multiplication in the file,
\texttt{mm\_a\_col\_major\_b\_col\_major}, is \textbf{not a GPU kernel at all}: it has no
\texttt{\_\_global\_\_} or \texttt{\_\_device\_\_} qualifier, takes no thread indices, and
is called directly from \texttt{main()} on the host. It is a plain, serial,
triple-nested-loop C++ function that runs on the \textbf{CPU}, used purely to compute a
correctness reference (\texttt{get\_avg\_abs\_diff\_ratio} compares it against the
tensor-core result) --- its runtime is never measured or reported, unlike the tensor-core
kernel's (\texttt{measure\_performance}, cudaEvent-timed).

\textbf{Summary: 100\% of GPU-side matrix multiplication in \texttt{leimao-mm.cu} is
tensor-core work. 0\% is SIMT-core.} A true SIMT GEMM kernel, as sketched in
Section~\ref{sec:fp32}, does not exist in \texttt{leimao-mm.cu} in any form --- it has
now been written from scratch as a separate pair of files,
\texttt{leimao-simt-variant-fp32.cu} and \texttt{leimao-simt-variant-fp16.cu}, benchmarked
directly against the tensor-core path in Section~\ref{sec:benchmark}.

\section{SIMT vs.\ Tensor Core: A Real Benchmark (FP16)}
\label{sec:benchmark}

\texttt{leimao-simt-variant-fp32.cu} and \texttt{leimao-simt-variant-fp16.cu} are the SIMT
kernels Section~\ref{sec:simt} found missing from \texttt{leimao-mm.cu} --- a standard
shared-memory tiled GEMM ($32\times32$ tiles, one thread per output element, ordinary
\texttt{FFMA} throughout, ordinary CUDA (SIMT) cores, no \texttt{wmma} namespace
anywhere), restricted to exactly the row-major $\times$ row-major case this report
studies (Section~\ref{sec:scope}). The FP16 variant matches
\texttt{leimao-mm.cu}'s own accumulation choice for a fair comparison: \texttt{\_\_half}
inputs, \texttt{float} accumulation (via \texttt{\_\_half2float} inside the inner loop),
\texttt{float} output --- the same precision contract as
\texttt{launch\_wmma\_mm<\_\_half, float>}, so the only thing actually being compared is
\emph{which hardware path performs the multiply-accumulate}, not a difference in
numerical precision along the way. Both SIMT files reuse the exact same
\texttt{measure\_performance} (cudaEvent-timed, 10 warmups, 10 measured repeats) and
$1024\times1024\times1024$ problem size as \texttt{leimao-mm.cu}, and both were checked
for correctness against a CPU reference before timing (no ``high average absolute diff
ratio'' warning printed by either --- i.e.\ within the same $0.01$ average relative-error
tolerance \texttt{leimao-mm.cu} itself uses).

\begin{center}
\begin{tabular}{lrr}
\toprule
& Latency (ms) & \\
\midrule
Tensor core (HMMA, \texttt{leimao-mm.cu}, ``C = A\^{}T * B\^{}T'') & 0.177--0.178 & (3 runs) \\
SIMT (\texttt{leimao-simt-variant-fp16.cu}) & 0.864--0.866 & (3 runs) \\
\bottomrule
\end{tabular}
\end{center}

Averaging each: tensor core $\approx 0.178$~ms, SIMT $\approx 0.865$~ms ---
\textbf{the tensor-core path is $\approx\!4.9\times$ faster} than this tiled SIMT kernel,
same GPU, same problem size, same input/accumulation precision, same $1024^3$-element
problem. This is a real, measured number from this machine, not a textbook or vendor
figure; it should not be read as ``the'' tensor-core speedup in general --- a more
aggressively tuned SIMT kernel (larger tiles, register blocking, double-buffered shared
memory, vectorized loads) would close some of this gap, and the gap itself grows or
shrinks with problem size and GPU generation. What it does establish concretely: even a
reasonably standard, textbook-tiled SIMT GEMM is not close to tensor-core throughput on
this hardware for this problem size --- the $16\times16\times16$-element HMMA instruction
Section~\ref{sec:fp16} walks through is doing, per issue, what this SIMT kernel needs
$16^3 = 4096$ separate scalar \texttt{FFMA}-equivalent operations to accomplish.

\section{Summary Table}

\begin{center}
\begin{tabular}{lllll}
\toprule
Precision & Tensor core? & Fragment type & Accumulator & Tile ($M{\times}N{\times}K$) \\
\midrule
FP16 & Yes (HMMA) & \texttt{\_\_half} & \texttt{\_\_half} \emph{or} \texttt{float} & $16{\times}16{\times}16$ (+2 more) \\
BF16 & Yes & \texttt{\_\_nv\_bfloat16} & \texttt{float} only & $16{\times}16{\times}16$ (+2 more) \\
FP32 (TF32) & Yes & \texttt{precision::tf32} & \texttt{float} & $16{\times}16{\times}\mathbf{8}$ \\
FP32 (full) & No & --- (SIMT only) & \texttt{float} & N/A \\
\bottomrule
\end{tabular}
\end{center}

\begin{center}
\begin{tabular}{ll}
\toprule
Precision & Exercised in \texttt{tc\_matmult/}? \\
\midrule
FP16, tensor core & Yes --- \texttt{leimao-mm.cu}'s only working GEMM path (FP32 accumulator) \\
FP16, SIMT & Yes --- \texttt{leimao-simt-variant-fp16.cu} (Section~\ref{sec:benchmark}) \\
FP32, SIMT & Yes --- \texttt{leimao-simt-variant-fp32.cu} \\
BF16 & No --- supported by the header/hardware, absent from this directory \\
FP32 (TF32) & No --- supported by the header/hardware, absent from this directory \\
\bottomrule
\end{tabular}
\end{center}

\paragraph{A note beyond the three precisions in scope.} The same header also defines
\texttt{fragment<matrix\_a/\allowbreak matrix\_b, 8, 8, 4, double,
row\_major$|$col\_major>} --- double-precision (DMMA) tensor-core matrix multiplication
exists too, at a much smaller $8{\times}8{\times}4$ tile. Out of scope for this report,
but worth knowing it isn't only FP16/BF16/TF32/INT8 that get tensor-core paths on this
architecture family.

\section{Conclusion}

For row-major $A$ times row-major $B$ (Section~\ref{sec:scope}): FP16 works today,
unmodified, via \texttt{leimao-mm.cu}'s existing \texttt{is\_A\_transpose=true,
is\_B\_transpose=true} call. BF16 is fully supported by the same CUDA/hardware stack,
confirmed directly in \texttt{mma.hpp}, but needs new code --- a direct, same-shape port
of the FP16 path with the element type swapped and the accumulator forced to
\texttt{float}. FP32 has no true full-precision tensor-core path at all; the tensor-core
option is TF32 (an explicit, hardware-managed precision reduction on load, full-precision
accumulate, smaller $K$-tile), and genuine full-precision FP32 GEMM would require an
entirely separate, non-tensor-core SIMT kernel. Every GPU-side multiply-accumulate that
exists in \texttt{leimao-mm.cu} itself runs on tensor cores; its only non-tensor-core
matrix multiply is a CPU-side correctness check, not a GPU kernel.

That missing SIMT kernel now exists, in both precisions (Section~\ref{sec:benchmark}),
and the FP16
comparison in
Section~\ref{sec:benchmark} puts a real number on the gap this report otherwise only
described structurally: on this GPU, at $1024\times1024\times1024$, with matched
input/accumulation precision, tensor cores finish in $\approx 0.178$~ms against a
reasonably standard tiled SIMT kernel's $\approx 0.865$~ms --- roughly $4.9\times$.
Section~\ref{sec:fp16}'s structural point (one HMMA instruction does the work of 4096
scalar multiply-adds) and this measured $4.9\times$ are two different kinds of evidence
for the same conclusion, not the same number restated: the instruction count says why a
gap should exist; the benchmark says how large it actually is for this kernel, this
problem size, this GPU.

\end{document}
